\section{Related Work}

{\bf Autonomous Scientific Discovery.} The development of AI agents capable of full research lifecycles has rapidly accelerated. \citet{lu2024aiscientist} introduced the first comprehensive framework for automated scientific discovery, handling everything from idea generation to paper writing. Building on this, \citet{yamada2025aiscientistv2} proposed an agentic tree search algorithm with vision-language model feedback to eliminate reliance on human-authored code templates. While these systems demonstrate the feasibility of automated science, they largely operate on known machine learning tasks where standard templates can be remixed. In contrast, our \methodname specifically evaluates whether these systems can step outside their training distribution to discover entirely novel, counterfactual relationships.

{\bf Benchmarking Scientific Reasoning.} Evaluating the scientific reasoning capabilities of LLMs requires moving beyond static QA datasets. \citet{chen2025physgym} introduced PhysGym to benchmark LLMs in interactive physics environments with controlled prior knowledge, demonstrating the importance of anonymizing variables to test for memorization. Similarly, \citet{zheng2025newtonbench} addressed the evaluation trilemma using counterfactual law shifts, systematically altering canonical laws to prevent reliance on memorized textbook formulas. \citet{koblischke2025gravitybench} focused on strategic observation planning in out-of-distribution gravitational scenarios. Building on these approaches, we operationalize counterfactual shifts within an end-to-end AI Scientist pipeline, evaluating not just the mathematical discovery but the resulting scientific manuscript.

{\bf Pitfalls of AI-Driven Science.} As AI systems take on more autonomous roles in research, understanding their failure modes becomes critical. \citet{luo2025hidden} provided a critical analysis of hidden pitfalls in AI Scientist systems, including post-hoc selection bias and metric misuse. Our findings complement this work by demonstrating another severe pitfall: the phenomenon of semantic override, where strong linguistic priors cause an AI agent to completely ignore conflicting empirical evidence and fabricate statistical validation in a generated research paper.